\begin{proof}
First observe that it follows from
\ref{sga3:XXI:1.2.10} and \ref{sga3:XXI:3.3.6} that (i) is a
consequence of (ii) and \SGARefLink{sga3:XXI:3.6:item:iii}{(iii)}, and that (ii) and \SGARefLink{sga3:XXI:3.6:item:iii}{(iii)} determine the
\(T_\alpha\) completely. We construct them by induction on
\(\operatorname{ord}_\Delta(\alpha)\). If
\(\operatorname{ord}_\Delta(\alpha)=1\), then \(\alpha\in\Delta\),
and we put \(T_\alpha=T_i\) if \(\alpha=\alpha_i\).

Consider the hypothesis
\[
\begin{split}
(\mathrm{H}_p)\quad&
\text{there exist \(T_\alpha\), for
\(\alpha\in\mathscr{P}(\Delta)\) with
\(\operatorname{ord}_\Delta(\alpha)\le p\), satisfying (ii),}\\
&\text{and satisfying (iii) whenever
\(\operatorname{ord}_\Delta(\alpha)\le p\) and
\(\operatorname{ord}_\Delta(\beta)\le p\).}
\end{split}
\]
It holds for \(p=1\). Indeed, if \(\alpha\) and
\(s_{\alpha_i}(\alpha)=\beta\) are simple, then \(\alpha_i\) and
\(\alpha\) are orthogonal. Thus, if
\(\alpha_j=\alpha=\beta\), one has \(n_{ij}=2\), whence
\[
 T_iT_jT_i=T_j.
\]

Suppose \(p>1\) and \((\mathrm{H}_{p-1})\) holds.

\medskip
\noindent\textbf{A) Construction of the \(T_\alpha\) for
\(\operatorname{ord}_\Delta(\alpha)\le p\).}
It is enough to do this when
\(\operatorname{ord}_\Delta(\alpha)=p\). There exists
\(\alpha_i\in\Delta\) such that
\[
 s_{\alpha_i}(\alpha)\in\mathscr{P}(\Delta),
 \qquad
 \operatorname{ord}_\Delta(s_{\alpha_i}(\alpha))<p
\]
by Lemma~\ref{sga3:XXI:3.3.3}. Put
\begin{equation}
\tag{\(\star\)}
\label{sga3:XXI:5.2:construction}
 T_\alpha=T_iT_{s_{\alpha_i}(\alpha)}T_i.
\end{equation}
We verify that \(T_\alpha\) depends only on \(\alpha\). Let
\(\alpha_j\in\Delta\) also satisfy
\[
 s_{\alpha_j}(\alpha)\in\mathscr{P}(\Delta),
 \qquad
 \operatorname{ord}_\Delta(s_{\alpha_j}(\alpha))<p.
\]
We must prove
\begin{equation}
\tag{\(+\)}
\label{sga3:XXI:5.2:independence}
 T_iT_{s_{\alpha_i}(\alpha)}T_i
 =T_jT_{s_{\alpha_j}(\alpha)}T_j.
\end{equation}
We distinguish two cases.

\begin{enumerate}
\item[(1)] Suppose that \(\alpha\) is a linear combination of
\(\alpha_i\) and \(\alpha_j\). The same is then true of
\(s_{\alpha_i}(\alpha)\) and \(s_{\alpha_j}(\alpha)\); by
\((\mathrm{H}_{p-1})\), the elements
\(T_{s_{\alpha_i}(\alpha)}\) and
\(T_{s_{\alpha_j}(\alpha)}\) are words in \(T_i,T_j\). Since the
projection of \eqref{sga3:XXI:5.2:independence} to \(W\) holds and
the theorem is true for \(n=2\), the restriction of \(p\) to the
subgroup of \(\widetilde W\) generated by \(T_i,T_j\) is injective.
Thus \eqref{sga3:XXI:5.2:independence} holds.

\item[(2)] Suppose that \(\alpha\) is not a linear combination of
\(\alpha_i\) and \(\alpha_j\). Then, for
\(w\in W_{\alpha_i,\alpha_j}\), all the \(w(\alpha)\) are positive
(cf. Lemma~\ref{sga3:XXI:3.4.9}). The relation to be verified can
also be written
\begin{equation}
\tag{\({+}{+}\)}
\label{sga3:XXI:5.2:rank-two-relation}
 (T_iT_j)^{n_{ij}-1}T_{s_{\alpha_i}(\alpha)}
 (T_jT_i)^{n_{ij}-1}
 =T_{s_{\alpha_j}(\alpha)}.
\end{equation}
Corollary~\ref{sga3:XXI:4.4} shows that all the \(w(\alpha)\) have
order \(<p\) for
\(w\in W_{\alpha_i,\alpha_j}\), \(w\ne1\). We may therefore apply
\((\mathrm{H}_{p-1})\) \(2(n_{ij}-1)\) times, which proves the
relation.
\end{enumerate}

\medskip
\noindent\textbf{B) Verification of \((\mathrm{H}_p)\).}
\footnote{Editors' note: The argument in the original has been
expanded in what follows.}
We must verify that, if \(\alpha_j\in\Delta\) and
\(\beta=s_{\alpha_j}(\alpha)\) satisfies
\(\operatorname{ord}_\Delta(\beta)\le p\), then
\[
 T_jT_\alpha T_j=T_\beta.
\]
If \(\operatorname{ord}_\Delta(\beta)<p\), this follows from the
preceding construction, since \(T_j^2=1\). We may therefore suppose
\[
 \operatorname{ord}_\Delta(\beta)
 =p=\operatorname{ord}_\Delta(\alpha).
\]
In this case \(\alpha\) and \(\alpha_j\) are orthogonal, so
\(\beta=\alpha\), and it remains to prove
\begin{equation}
\tag{\(\dagger\)}
\label{sga3:XXI:5.2:orthogonal}
 T_jT_\alpha T_j=T_\alpha.
\end{equation}
By \eqref{sga3:XXI:5.2:construction},
\[
 T_\alpha=T_iT_{s_{\alpha_i}(\alpha)}T_i,
\]
so it remains to verify
\begin{equation}
\tag{\({+}{+}{+}\)}
\label{sga3:XXI:5.2:orthogonal-expanded}
 T_jT_iT_{s_{\alpha_i}(\alpha)}T_iT_j
 =T_\alpha
 =T_iT_{s_{\alpha_i}(\alpha)}T_i.
\end{equation}
\end{proof}
